\section{Finite Galois, profinite, and Tannakian resolved symmetry}
\label{sec:block12-galois-tannakian}

Galois-Tannakian atlas develops the final declared stratum in the present eight-stratum program: finite semisimple Galois charts, cohomological descent, profinite refinement, and the tensor-fiber boundary at which a symmetry group can be reconstructed.  The strongest Galois source already contains Reynolds projectors, the subgroup--fixed-field correspondence, normality and quotient compression, conjugate projector atlases, Mackey double-coset packets, ramification details, Haar projectors, Hilbert 90, the Brauer obstruction, finite \(\acute{e}\)tale permutation resolutions, and neutral Tannakian reconstruction \cite{GAL}.  The purpose here is to place those results into the same passport, bridge, obstruction, and effective-descent language used by the preceding sections.

The scope is deliberately semisimple or tensor-exact.  For finite group averaging, the coefficient field has characteristic zero or characteristic prime to the relevant group order.  Linear semilinear descent is taken over a finite Galois extension.  Projective descent is typed separately from vector descent.  Tannakian reconstruction is asserted only after a faithful exact tensor fiber functor has been supplied.  Modular fixed points, wild integral descent, nonneutral fiber functors, and derived or flat replacements for failed averaging remain outside this block.

\subsection{Finite subgroup, fixed-field, and Reynolds-projector charts}
\label{subsec:block12-reynolds}

Let \(L/K\) be a finite Galois extension with group
\[
 G=\operatorname{Gal}(L/K).
\]
Let \(V\) be a finite-dimensional representation of \(G\) over a characteristic-zero field, or a unitary complex realization obtained after scalar extension.  For every subgroup \(H\leq G\), define
\[
 \boxed{
 P_H=\frac1{|H|}\sum_{h\in H}U_h.}
 \tag{12.1}
\]

\begin{theorem}[Reynolds fixed packet and reverse order]
\label{thm:block12-reynolds}
The operator \(P_H\) is the projector onto the fixed packet
\[
 V^H=\{v:U_hv=v\text{ for all }h\in H\}.
\]
If \(H\leq J\), then
\[
 \boxed{P_HP_J=P_JP_H=P_J.}
 \tag{12.2}
\]
Consequently, subgroup inclusion reverses the fixed-projector order.
\end{theorem}

\begin{proof}
The average is invariant under left and right multiplication by elements of \(H\).  Therefore \(P_H^2=P_H\), and its range is exactly the fixed subspace.  If \(H\leq J\), then \(V^J\subseteq V^H\).  The two projections are nested and multiply to the projector with smaller range.
\end{proof}

In the faithful regular realization, or in the field realization supplied by a normal basis, the projector determines the subgroup.  Indeed, the support of the group-algebra idempotent
\[
 e_H=\frac1{|H|}\sum_{h\in H}h
\]
is exactly \(H\).

\begin{theorem}[Normalized finite Galois chart]
\label{thm:block12-normalized-galois-chart}
Fix the finite Galois extension \(L/K\).  The following data determine one another exactly:
\[
 \boxed{
 H
 \quad\longleftrightarrow\quad
 L^H
 \quad\longleftrightarrow\quad
 P_H
 }
 \tag{12.3}
\]
where the projector is taken in a faithful field or regular realization.  The correspondence is order reversing:
\[
 H\leq J
 \quad\Longleftrightarrow\quad
 L^H\supseteq L^J
 \quad\Longleftrightarrow\quad
 P_H\geq P_J.
\]
Thus the subgroup poset, the intermediate-field poset, and the faithful fixed-projector poset form three normalized charts of one finite semisimple Galois object.
\end{theorem}

\begin{proof}
The finite Galois correspondence gives \(H=\operatorname{Gal}(L/L^H)\) and \(E=L^{\operatorname{Gal}(L/E)}\).  The range of \(P_H\) is \(L^H\) after faithful scalar extension, and the support of \(e_H\) recovers \(H\).  All three assignments are therefore mutually inverse on the declared charts.
\end{proof}

\begin{figure}[htbp]
\centering
\includegraphics[width=.97\textwidth]{galois_subgroup_atlas.pdf}
\caption{The finite subgroup, Reynolds-projector, and fixed-field charts are reversible in a faithful semisimple realization.  Normality is the additional condition required for quotient compression.}
\label{fig:block12-subgroup-atlas}
\end{figure}

\subsection{Compatibility, normality, and conjugate atlases}
\label{subsec:block12-normality}

For two subgroups \(H,J\leq G\), the common fixed packet is always
\[
 V^H\cap V^J=V^{\langle H,J\rangle}.
\]
The two averaging projectors need not commute.

\begin{theorem}[Permutable subgroups and exact common refinement]
\label{thm:block12-permutable}
In a faithful regular realization, the following are equivalent:
\begin{enumerate}
\item \(HJ=JH\);
\item \(HJ\) is a subgroup of \(G\);
\item \(P_HP_J=P_JP_H\).
\end{enumerate}
When these conditions hold,
\[
 \boxed{P_HP_J=P_{HJ}.}
 \tag{12.4}
\]
If they fail, \(P_HP_J\) and \(P_JP_H\) are mixed atlas corners rather than projectors.
\end{theorem}

\begin{proof}
If \(HJ=JH\), every element of \(HJ\) has \(|H\cap J|\) representations as \(hj\).  Multiplying the two group averages therefore gives the average over \(HJ\).  Conversely, the support of \(e_He_J\) is \(HJ\), while the support of \(e_Je_H\) is \(JH\).  Equality in the regular group algebra forces equality of the product sets.
\end{proof}

\begin{theorem}[Normality and quotient compression]
\label{thm:block12-normality}
For a subgroup \(H\leq G\), the following conditions are equivalent:
\begin{enumerate}
\item \(H\triangleleft G\);
\item \(e_H\) is central in the group algebra;
\item \(P_HU_g=U_gP_H\) for every \(g\in G\);
\item the fixed packet \(V^H\) is \(G\)-invariant;
\item the fixed field \(L^H\) is Galois over \(K\).
\end{enumerate}
When they hold, the compressed action on \(V^H\) factors through \(G/H\), and
\[
 \boxed{\operatorname{Gal}(L^H/K)\cong G/H.}
 \tag{12.5}
\]
\end{theorem}

\begin{proof}
Conjugation sends \(e_H\) to \(e_{gHg^{-1}}\).  Centrality is therefore equivalent to normality.  The represented identity
\[
 U_gP_HU_g^{-1}=P_{gHg^{-1}}
\]
gives the operator conditions.  The field and quotient statements are the normal-subgroup part of the finite Galois correspondence.
\end{proof}

If \(H\) is not normal, its conjugates form a genuine atlas rather than one quotient chart.  Put
\[
 \operatorname{Core}_G(H)=\bigcap_{g\in G}gHg^{-1},
 \qquad
 \operatorname{Ncl}_G(H)=\left\langle gHg^{-1}:g\in G\right\rangle.
\]

\begin{theorem}[Canonical envelopes of a nonnormal Galois atlas]
\label{thm:block12-core-normal-closure}
Let \(M=L^H\).  Then
\[
 \boxed{
 L^{\operatorname{Core}_G(H)}
 =\prod_{g\in G}g(M)}
 \tag{12.6}
\]
is the smallest Galois extension inside \(L\) containing every conjugate chart, while
\[
 \boxed{
 L^{\operatorname{Ncl}_G(H)}
 =\bigcap_{g\in G}g(M)}
 \tag{12.7}
\]
is the largest Galois subfield contained in every chart.  The core and normal-closure projectors are the two canonical global envelopes of the conjugate atlas.
\end{theorem}

\begin{proof}
Composita of intermediate fields correspond to intersections of subgroups, while intersections of fields correspond to generated subgroups.  The core is the largest normal subgroup inside \(H\), and the normal closure is the smallest normal subgroup containing \(H\).
\end{proof}

\begin{example}[The \(S_3\) nonnormal atlas]
Let \(G=S_3\) and let \(H\) be generated by one transposition.  There are three conjugate order-two subgroups and hence three conjugate cubic fixed-field charts.  Their subgroup core is trivial, while their normal closure is all of \(S_3\).  By contrast, the order-three subgroup \(A_3\) is normal and its fixed chart carries the quotient symmetry \(S_3/A_3\cong C_2\).
\end{example}

\subsection{Double cosets as exact transition packets}
\label{subsec:block12-mackey}

Restriction and induction connect subgroup charts that need not commute.  Their exact address set is the double-coset space
\[
 K\backslash G/H.
\]

\begin{theorem}[Mackey transition-packet decomposition]
\label{thm:block12-mackey}
Let \(H,K\leq G\), let \(W\) be a finite-dimensional representation of \(H\), and choose representatives \(x\) for \(K\backslash G/H\).  Then
\[
 \boxed{
 \operatorname{Res}_K^G\operatorname{Ind}_H^GW
 \cong
 \bigoplus_{x\in K\backslash G/H}
 \operatorname{Ind}_{K\cap xHx^{-1}}^K
 \left({}^x\operatorname{Res}_{H\cap x^{-1}Kx}^HW\right).}
 \tag{12.8}
\]
Each double coset is one exact transfer packet between the \(H\)- and \(K\)-charts.
\end{theorem}

\begin{proof}
The induced representation is a function space on \(G\) with \(H\)-covariance.  Functions supported on different double cosets are invariant direct summands under \(K\).  Identifying the summand on \(KxH\) with an induced representation from \(K\cap xHx^{-1}\) gives the decomposition.
\end{proof}

The dimension identity for the trivial \(H\)-representation is
\[
 \boxed{
 [G:H]
 =
 \sum_{x\in K\backslash G/H}
 [K:K\cap xHx^{-1}].}
 \tag{12.9}
\]
This is the finite packet-count law audited by the Galois-Tannakian atlas verifier.

\begin{figure}[htbp]
\centering
\includegraphics[width=.97\textwidth]{galois_mackey_packets.pdf}
\caption{A nonnormal fixed object is an atlas of conjugate charts.  Core and normal closure give two global envelopes, while double cosets give exact mixed transition packets.}
\label{fig:block12-mackey}
\end{figure}

\begin{table}[htbp]
\centering
\caption{Finite group models used in the Galois-Tannakian atlas subgroup-projector audit.}
\label{tab:block12-groups}
\scriptsize
\begin{tabularx}{\textwidth}{@{}L{0.12\textwidth}rrrrrY@{}}
\toprule
Group & Order & Subgroups & Normal & Permutable ordered pairs & Noncommuting ordered pairs & Double-coset packets audited\\
\midrule
C2 & 2 & 2 & 2 & 4 & 0 & 5\\
C3 & 3 & 2 & 2 & 4 & 0 & 6\\
C4 & 4 & 3 & 3 & 9 & 0 & 15\\
C5 & 5 & 2 & 2 & 4 & 0 & 8\\
C6 & 6 & 4 & 4 & 16 & 0 & 30\\
C8 & 8 & 4 & 4 & 16 & 0 & 37\\
V4 & 4 & 5 & 5 & 25 & 0 & 37\\
S3 & 6 & 6 & 3 & 30 & 6 & 65\\
D4 & 8 & 10 & 6 & 92 & 8 & 195\\
A4 & 12 & 10 & 3 & 64 & 36 & 250\\
\bottomrule
\end{tabularx}
\end{table}


\subsection{Ramification as a nested resolution}
\label{subsec:block12-ramification}

Let \(L/K\) be a finite Galois extension of local fields.  Write
\[
 G_0\supseteq G_1\supseteq G_2\supseteq\cdots
\]
for the lower ramification filtration after inertia, and let \(V\) be a finite-dimensional complex representation of \(G\).  Define
\[
 P_i=\frac1{|G_i|}\sum_{g\in G_i}U_g.
\]
The ramification groups are nested and normal, so
\[
 P_0\leq P_1\leq P_2\leq\cdots\leq I.
\]
Put
\[
 D_i=P_{i+1}-P_i.
\]

\begin{theorem}[Orthogonal ramification details]
\label{thm:block12-ramification-details}
The operators \(D_i\) are mutually orthogonal projectors, and
\[
 \boxed{
 V=V^{G_0}\oplus\bigoplus_i D_iV.}
 \tag{12.10}
\]
Every detail packet is stable under the quotient action permitted by the corresponding normal subgroup.
\end{theorem}

\begin{proof}
Nested orthogonal projectors commute, and the difference of two nested projectors is a projector onto the orthogonal difference of their ranges.  The decomposition follows by telescoping.
\end{proof}

\begin{theorem}[Artin conductor as weighted unresolved rank]
\label{thm:block12-artin-conductor}
The Artin and Swan conductors satisfy
\[
 \boxed{
 a(V)=\sum_{i\geq0}\frac{|G_i|}{|G_0|}\operatorname{Tr}(I-P_i),}
 \tag{12.11}
\]
\[
 \boxed{
 \operatorname{Sw}(V)=\sum_{i\geq1}\frac{|G_i|}{|G_0|}\operatorname{Tr}(I-P_i).}
 \tag{12.12}
\]
If \(d_j=\operatorname{rank}D_j\), then
\[
 \boxed{
 a(V)=\sum_jd_j\sum_{i=0}^{j}\frac{|G_i|}{|G_0|}.}
 \tag{12.13}
\]
Thus the conductor is the weighted total rank that remains unresolved at each ramification scale.
\end{theorem}

\begin{proof}
The classical conductor formula uses \(\operatorname{codim}V^{G_i}\).  Since \(P_i\) projects onto \(V^{G_i}\), this codimension is \(\operatorname{Tr}(I-P_i)\).  Also
\[
 \operatorname{codim}V^{G_i}=\sum_{j\geq i}d_j.
\]
Reversing the finite sums gives the detail formula.
\end{proof}

\begin{figure}[htbp]
\centering
\includegraphics[width=.92\textwidth]{galois_ramification_resolution.pdf}
\caption{Ramification subgroups give nested fixed projectors.  Their orthogonal differences are the resolved tame and wild detail packets, and the conductor is their weighted unresolved rank.}
\label{fig:block12-ramification}
\end{figure}

\subsection{Profinite completion and finite quotient effectivity}
\label{subsec:block12-profinite}

Let \(G\) be a compact group and let \(U:G\to\mathcal U(\mathcal H)\) be strongly continuous.  For a closed subgroup \(H\), normalized Haar measure defines
\[
 \boxed{
 P_H=\int_HU_h\,d\mu_H(h).}
 \tag{12.14}
\]
The same averaging proof as in the finite case shows that \(P_H\) is the orthogonal projector onto \(\mathcal H^H\).

\begin{theorem}[Shrinking-subgroup reconstruction]
\label{thm:block12-shrinking-subgroups}
Let \(G\) be profinite and let
\[
 U_1\supseteq U_2\supseteq\cdots
\]
be a neighborhood basis of the identity consisting of open normal subgroups.  Then
\[
 \boxed{P_{U_n}\longrightarrow I}
 \tag{12.15}
\]
strongly in every strongly continuous unitary representation.  More generally, if \(N\triangleleft G\) is closed and \(U_n/N\) shrinks to the identity in \(G/N\), then \(P_{U_n}\to P_N\) strongly.
\end{theorem}

\begin{proof}
For each \(v\in\mathcal H\), strong continuity gives
\[
 \sup_{u\in U_n}\|U_uv-v\|\longrightarrow0.
\]
Averaging this estimate over \(U_n\) proves the first statement.  Apply it to the quotient action on the \(N\)-fixed packet for the second.
\end{proof}

\begin{corollary}[Finite-dimensional representations occur at finite level]
\label{cor:block12-finite-level}
Every continuous finite-dimensional complex representation of a profinite group has finite image and factors through a finite quotient.
\end{corollary}

\begin{proof}
The image is a compact totally disconnected subgroup of \(\operatorname{GL}_n(\mathbb C)\).  After conjugation it lies in a compact Lie group.  A compact totally disconnected Lie group is finite.
\end{proof}

This finite-level factorization is a strong effectivity statement for finite-dimensional complex representations.  It is not a statement that every integral, modular, or derived representation becomes semisimple at finite level.

\subsection{Semilinear vector descent and Hilbert 90}
\label{subsec:block12-hilbert90}

Let \(L/K\) be finite Galois with group \(G\).  A semilinear \(G\)-object is an \(L\)-vector space \(V\) with maps
\[
 \rho_\sigma:V\to V,
 \qquad
 \rho_\sigma(\lambda v)=\sigma(\lambda)\rho_\sigma(v),
\]
satisfying the group law.  Scalar extension and fixed vectors define functors
\[
 L\otimes_K-:\operatorname{Vect}_K\longrightarrow\operatorname{Semilin}_G(L),
\]
\[
 (-)^G:\operatorname{Semilin}_G(L)\longrightarrow\operatorname{Vect}_K.
\]

\begin{theorem}[Effective finite Galois vector descent]
\label{thm:block12-vector-descent}
The two functors above are quasi-inverse equivalences:
\[
 \boxed{
 \operatorname{Vect}_K
 \simeq
 \operatorname{Semilin}_G(L).}
 \tag{12.16}
\]
For every semilinear object \(V\), the fixed packet \(W=V^G\) satisfies
\[
 \boxed{W\otimes_KL\cong V.}
 \tag{12.17}
\]
The equivalence is effective on both objects and morphisms.
\end{theorem}

\begin{proof}
This is finite Galois descent.  After an \(L\)-basis is chosen, the semilinear action is represented by a cocycle
\[
 a_{\sigma\tau}=a_\sigma\sigma(a_\tau),
 \qquad
 a_\sigma\in\operatorname{GL}_n(L).
\]
Nonabelian Hilbert 90 gives \(B\in\operatorname{GL}_n(L)\) with
\[
 a_\sigma=B^{-1}\sigma(B).
\]
The columns of \(B^{-1}\) applied to a \(K\)-basis are fixed by the twisted action.  They form an \(n\)-dimensional fixed packet, and scalar extension reconstructs \(V\).  An equivariant \(L\)-linear map restricts to a \(K\)-linear map on fixed packets, and scalar extension recovers it.
\end{proof}

A useful constructive form is obtained by choosing any matrix \(C\in M_n(L)\) and setting
\[
 T(C)=\sum_{\sigma\in G}a_\sigma\sigma(C).
\]
Then
\[
 a_\tau\tau(T(C))=T(C).
\]
Whenever \(T(C)\) is invertible,
\[
 B=T(C)^{-1}
\]
is a Hilbert 90 trivialization.  Standard descent guarantees that such a \(C\) exists.

\begin{definition}[Framed semilinear descent chart]
A framed descent chart retains an ordered \(K\)-basis of the fixed packet \(W\).  Its scalar extension gives a distinguished \(L\)-basis of \(V\), and hence a literal cocycle matrix family.
\end{definition}

\begin{proposition}[Gauge and intrinsic automorphisms remain distinct]
\label{prop:block12-framed-unframed}
For fixed rank \(n\), nonabelian Hilbert 90 implies that every \(\operatorname{GL}_n(L)\)-cocycle is gauge equivalent to the trivial cocycle.  The unframed descent groupoid is not contractible: its automorphism group is \(\operatorname{GL}_n(K)\).  After a global \(K\)-frame is retained, the normalized fiber becomes contractible.
\end{proposition}

\begin{proof}
Hilbert 90 gives one isomorphism class.  Automorphisms of the descended \(K\)-space are exactly \(\operatorname{GL}_n(K)\).  A chosen ordered basis kills those automorphisms.
\end{proof}

\subsection{Projective descent and the Brauer obstruction}
\label{subsec:block12-brauer}

A projective cocycle takes values in \(\operatorname{PGL}_n(L)\).  Choose lifts
\[
 A_\sigma\in\operatorname{GL}_n(L).
\]
Their cocycle failure is scalar:
\[
 \boxed{
 c(\sigma,\tau)
 =A_\sigma\sigma(A_\tau)A_{\sigma\tau}^{-1}
 \in L^\times.}
 \tag{12.18}
\]
Changing the lifts alters \(c\) by a scalar coboundary.

\begin{theorem}[Exact projective descent obstruction]
\label{thm:block12-brauer}
The scalar class
\[
 \boxed{
 [c]=\delta([\alpha])
 \in H^2(G,L^\times)
 =\operatorname{Br}(L/K)}
 \tag{12.19}
\]
is independent of the chosen lifts.  It vanishes exactly when the projective cocycle lifts to a linear semilinear cocycle.  In that case \cref{thm:block12-vector-descent} produces a descended vector space.  If it does not vanish, the descended object is a module over a central simple \(K\)-algebra split by \(L\), not an ordinary global vector packet.
\end{theorem}

\begin{proof}
Apply the connecting map in nonabelian cohomology to
\[
 1\longrightarrow L^\times
 \longrightarrow\operatorname{GL}_n(L)
 \longrightarrow\operatorname{PGL}_n(L)
 \longrightarrow1.
\]
Vanishing is equivalent to rescaling the lifts so that the scalar defect is removed.  The lifted cocycle then descends by Hilbert 90.  The nonzero classes are the relative Brauer classes of central simple algebras split by \(L\).
\end{proof}

\begin{example}[The quaternionic obstruction over \(\mathbb C/\mathbb R\)]
Let \(G=\{1,\sigma\}\), where \(\sigma\) is complex conjugation, and put
\[
 A_\sigma=
 \begin{pmatrix}
 0&1\\
 -1&0
 \end{pmatrix}.
\]
Then
\[
 A_\sigma\sigma(A_\sigma)=-I.
\]
The projective cocycle condition holds because \(-I\) is scalar.  Its boundary is represented by \(-1\) in
\[
 \mathbb R^\times/N_{\mathbb C/\mathbb R}(\mathbb C^\times).
\]
Every complex norm is positive, so \(-1\) is nontrivial.  The descended central simple algebra is the Hamilton quaternion algebra.  Replacing \(A_\sigma\) by
\[
 \begin{pmatrix}0&1\\1&0\end{pmatrix}
\]
gives scalar defect \(+1\) and hence split vector descent.
\end{example}

\begin{figure}[htbp]
\centering
\includegraphics[width=.97\textwidth]{galois_descent_obstruction.pdf}
\caption{Linear cocycles are removed by Hilbert 90.  Projective cocycles can satisfy every visible \(\operatorname{PGL}\) overlap law while retaining a nonzero degree-two Brauer obstruction.}
\label{fig:block12-descent-obstruction}
\end{figure}

\subsection{Finite \texorpdfstring{\(\acute{e}\)}{e}tale and neutral Tannakian symmetry}
\label{subsec:block12-tannakian}

For a connected scheme \(X\) with geometric point \(\bar x\), the fiber functor on finite \(\acute{e}\)tale covers identifies that category with finite continuous \(\pi_1^{\acute{e}t}(X,\bar x)\)-sets.  Every finite fiber \(S\) produces a permutation resolution on \(\ell^2(S)\): coordinate projectors are indexed by points, orbit projectors by transitive components, and covering morphisms give equivariant packet transfers.  The profinite fundamental group is the automorphism group of this fiber functor.

The tensor-linear version requires more structure.

\begin{theorem}[Neutral Tannakian reconstruction]
\label{thm:block12-tannakian}
Let \(\mathcal C\) be a neutral Tannakian category over \(k\), and let
\[
 \omega:\mathcal C\longrightarrow\operatorname{Vect}_k
\]
be a faithful exact tensor fiber functor.  Then
\[
 \boxed{
 \mathcal G=\operatorname{Aut}^{\otimes}(\omega)}
 \tag{12.20}
\]
is an affine group scheme, and
\[
 \boxed{
 \mathcal C\simeq\operatorname{Rep}_k(\mathcal G)}
 \tag{12.21}
\]
as tensor categories.  Declared packet idempotents and refinement maps preserved by \(\omega\) become tensor-compatible concrete projector resolutions in every fiber.
\end{theorem}

\begin{proof}
The group scheme and tensor equivalence are the neutral Tannakian reconstruction theorem.  Additivity and exact tensor preservation transport the declared idempotent and refinement equations to the fibers.
\end{proof}

\begin{boundary}[Faithful corner realization is not a fiber functor]
A faithful additive or Hilbert-space realization can preserve packet idempotents, transfer blocks, and operator equations without preserving tensor products, duals, exact sequences, or the tensor unit.  Such a realization does not by itself determine \(\operatorname{Aut}^{\otimes}(\omega)\).  Tannakian symmetry begins only after the faithful exact tensor gate has been verified.
\end{boundary}

\begin{proposition}[Finite group model of the tensor boundary]
\label{prop:block12-finite-tensor-model}
Let \(G\) be finite and let \(\omega:\operatorname{Rep}_k(G)\to\operatorname{Vect}_k\) be the forgetful functor.  The regular representation is faithful, the action on tensor products is diagonal, the unit is fixed, and the dual action is inverse transpose.  These data distinguish every \(g\in G\) and realize the finite constant-group case of \cref{thm:block12-tannakian}.
\end{proposition}

\begin{proof}
Left translation on the regular representation is faithful.  Functoriality of the tensor product, unit, and dual actions is immediate from the representation axioms.  Since different group elements act differently on the regular basis, the induced tensor automorphisms are distinct.
\end{proof}

\begin{figure}[htbp]
\centering
\includegraphics[width=.97\textwidth]{galois_profinite_tannakian.pdf}
\caption{Finite quotient projectors converge strongly in profinite realizations.  The lower row isolates the additional exact tensor structure required before a symmetry group scheme can be reconstructed.}
\label{fig:block12-profinite-tannakian}
\end{figure}

\subsection{The assembled finite Galois resolved atlas}
\label{subsec:block12-assembled}

The preceding results can be collected into one typed finite atlas.  The reversible charts are:
\begin{enumerate}
\item subgroups \(H\leq G\), fixed fields \(L^H\), and faithful Reynolds projectors \(P_H\);
\item finite-dimensional \(K\)-vector spaces and semilinear \(L\)-spaces with effective descent data;
\item finite \(\acute{e}\)tale fibers and their continuous permutation actions;
\item neutral Tannakian categories equipped with a specified faithful exact tensor fiber functor.
\end{enumerate}
Directional or obstructed bridges include:
\begin{enumerate}
\item restriction and induction between subgroup charts, decomposed by Mackey packets;
\item passage from a nonnormal chart to its core or normal-closure envelope;
\item projectivization from linear to projective cocycles;
\item forgetting tensor structure from a Tannakian fiber functor to an additive or corner realization;
\item passage from finite quotients to a profinite inverse system.
\end{enumerate}

\begin{theorem}[Finite semisimple Galois and Tannakian assembly]
\label{thm:block12-galois-assembly}
Under the hypotheses printed in this section:
\begin{enumerate}
\item the subgroup, fixed-field, and faithful fixed-projector charts form an order-reversing normalized atlas;
\item normal charts admit quotient compression, while nonnormal charts form conjugate projector atlases with canonical core and normal-closure envelopes;
\item double cosets give exact transition packets for induction and restriction;
\item ramification filtrations give nested fixed projectors whose weighted unresolved ranks are the Artin and Swan conductors;
\item continuous finite-dimensional profinite representations are detected at a finite quotient level;
\item linear semilinear descent is effective on objects and morphisms;
\item projective descent is effective exactly when its relative Brauer class vanishes;
\item neutral Tannakian symmetry is reconstructed exactly when a faithful exact tensor fiber functor is supplied.
\end{enumerate}
These statements form a complete assembly theorem for the declared finite semisimple and neutral tensor strata.  They do not imply effective descent for modular, wild, nonneutral, or arbitrary heterogeneous atlases.
\end{theorem}

\begin{proof}
Items 1--5 are \cref{thm:block12-normalized-galois-chart,thm:block12-normality,thm:block12-core-normal-closure,thm:block12-mackey,thm:block12-artin-conductor,cor:block12-finite-level}.  Item 6 is \cref{thm:block12-vector-descent}.  Item 7 is \cref{thm:block12-brauer}.  Item 8 is \cref{thm:block12-tannakian}.  The final boundary follows because the hypotheses used by Reynolds averaging, Hilbert 90, and neutral Tannakian reconstruction are not automatic in the excluded settings.
\end{proof}

A complete chart record on the finite subgroup locus is
\[
 \mathcal J_{L/K}(H)
 =
 \left(
 H,
 L^H,
 P_H,
 \operatorname{Core}_G(H),
 \operatorname{Ncl}_G(H),
 \{K\backslash G/H\}_{K\leq G}
 \right).
\]
The first three entries already classify the normalized object, while the remaining entries record its atlas position and transition structure.  For a framed semilinear vector object, the descended \(K\)-space with its frame is complete.  For a projective object, the Brauer class is an obstruction coordinate, not a complete invariant of all central simple algebras without the remaining algebra data.

\subsection{Exact finite verification}
\label{subsec:block12-verification}

The Galois-Tannakian atlas verifier uses exact rational group-algebra arithmetic, exact finite permutation combinatorics, and exact Gaussian-rational matrix algebra.  It audits cyclic groups, \(V_4\), \(S_3\), \(D_4\), and \(A_4\); every subgroup Reynolds idempotent; subgroup inclusion and projector order; normality and centrality; core and normal closure; commuting and noncommuting subgroup pairs; double-coset partitions and the Mackey dimension law; cyclic ramification-style filtrations; explicit Hilbert 90 coboundaries over \(\mathbb C/\mathbb R\); the split and quaternionic projective examples; and faithful regular tensor actions.

\begin{table}[htbp]
\centering
\caption{Exact verification of the Galois-Tannakian atlas finite Galois and Tannakian models.}
\label{tab:block12-validation}
\small
\begin{tabularx}{\textwidth}{@{}Y r@{}}
\toprule
Validation class & Checks\\
\midrule
Finite subgroup and Reynolds-projector identities & 1,020\\
Normality, core, and normal-closure identities & 192\\
Double-coset and Mackey packet identities & 1,590\\
Nested ramification-style projector filtrations & 39\\
Exact Hilbert 90 cocycles over Gaussian rationals & 884\\
Projective and Brauer obstruction examples & 44\\
Finite faithful tensor-model checks & 7,426\\
\midrule
\textbf{Galois-Tannakian atlas exact checks} & \textbf{11,195}\\
\textbf{Failures} & \textbf{0}\\
\bottomrule
\end{tabularx}
\end{table}


The verifier records \(11{,}195\) exact checks and no failures.  Together with the preceding certified atlas program, the cumulative exact record is
\[
 \boxed{2{,}038{,}860}
\]
checks with no recorded failures.

\begin{boundarybox}
\textbf{Exact boundary of Galois-Tannakian atlas.}
The finite subgroup/projector identities, semilinear vector descent, relative Brauer obstruction, profinite finite-level resolution, and neutral Tannakian tensor gate are established under their printed hypotheses.  The block does not prove exactness of fixed points when the group order is not invertible, derived replacement of failed Reynolds projectors, integral wild descent, classification of all central simple algebras from the Brauer class alone, existence of a neutral fiber functor, or effective descent for the entire heterogeneous Komorebi atlas.
\end{boundarybox}

\begin{resultbox}
\textbf{Section result.}
Finite semisimple Galois theory now appears as a fully typed resolved stratum.  Subgroups, fixed fields, and faithful fixed projectors are reversible charts.  Nonnormality produces conjugate atlases and mixed Mackey transition packets.  Ramification is a nested detail resolution.  Hilbert 90 makes linear semilinear descent effective, while the relative Brauer class is the exact obstruction to projective vector descent.  Profinite quotient charts reconstruct continuous finite-dimensional symmetry, and a faithful exact tensor fiber functor crosses the separate Tannakian boundary to reconstruct an affine group scheme.  This closes the declared finite Galois and neutral tensor stratum, not modular, wild, nonneutral, or unrestricted global descent.
\end{resultbox}
